% Wersja źródłowa: audyt II, jednostka 11
\setcounter{section}{10}
\section{Przestrzenie Frécheta. Twierdzenia o odwzorowaniu otwartym i wykresie domkniętym}\label{sec:frechet-open-graph}

\subsection{Dlaczego zupełność i metryzowalność występują razem?}

W \hyperref[sec:beczkowe-bornologiczne]{rozdziale 10} udowodniliśmy twierdzenie
Baire'a dla zupełnych przestrzeni metrycznych i wynikającą z niego
beczkowość przestrzeni Frécheta. W tym rozdziale wyprowadzimy dwa dalsze
fundamentalne rezultaty: ciągły liniowy operator surjektywny między
przestrzeniami Frécheta jest otwarty, a operator liniowy określony na
całej przestrzeni Frécheta jest ciągły, jeśli jego wykres jest domknięty
w iloczynie dwóch przestrzeni Frécheta.

W całym rozdziale $\K\in\{\R,\C\}$. Wszystkie przestrzenie
liniowo-topologiczne są Hausdorffa, chyba że podano inaczej.
Nie utożsamiamy zupełności względem kanonicznej struktury jednostajnej
z wymaganiem zbieżności jedynie ciągów Cauchy'ego bez wcześniejszego
sprawdzenia metryzowalności.

\begin{definicja}[Przestrzeń Frécheta]\label{def:frechet11}
Przestrzeń liniowo-topologiczną $E$ nad $\K$ nazywamy
\emph{przestrzenią Frécheta}, jeżeli $E$ jest lokalnie wypukła,
Hausdorffa, metryzowalna i zupełna względem swej kanonicznej
struktury jednostajnej. Oznaczenie $\mathcal L(E,F)$ oznacza
przestrzeń ciągłych operatorów liniowych $E\to F$.
\end{definicja}

\begin{stwierdzenie}[Metryka wyznaczona przez seminormy]\label{prop:frechet-metric11}
Niech $E$ będzie lokalnie wypukłą TVS Hausdorffa, której topologię
wyznacza przeliczalna rodzina ciągłych seminorm $(p_n)_{n\ge1}$.
Wtedy wzór
\[
 d(x,y)=\sum_{n=1}^{\infty}2^{-n}\min\{1,p_n(x-y)\},\qquad x,y\in E,
\]
określa metrykę niezmienniczą względem przesunięć, generującą
początkową topologię $E$. Metryka $d$ jest zupełna wtedy i tylko
wtedy, gdy $E$ jest zupełna jako TVS.
\end{stwierdzenie}
\begin{proof}
Każdy składnik jest nieujemny i niezmienniczy względem przesunięć.
Z nierówności $\min\{1,a+b\}\le \min\{1,a\}+\min\{1,b\}$
wynika nierówność trójkąta. Ponieważ seminormy rozdzielają punkty,
$d(x,y)=0$ implikuje $x=y$.

Jeśli $x_i\to x$ we wszystkich seminormach, to dla dowolnego
$\varepsilon>0$ wybieramy $N$ z $\sum_{n>N}2^{-n}<\varepsilon/2$.
Następnie pierwsze $N$ składników sumy definiującej $d(x_i,x)$
są dostatecznie małe, więc $d(x_i,x)<\varepsilon$.
Odwrotnie, z $d(x_i,x)\to0$ wynika dla ustalonego $n$,
że $\min\{1,p_n(x_i-x)\}\to0$, a więc $p_n(x_i-x)\to0$.
To dowodzi zgodności topologii, również dla sieci.

Z tego samego oszacowania wynika równoważność warunku Cauchy'ego
w metryce $d$ i warunku Cauchy'ego dla każdej seminormy $p_n$.
Metryka $d$ wyznacza więc kanoniczną strukturę jednostajną TVS,
a jej zupełność jest równoważna zupełności przestrzeni $E$.
\end{proof}

\begin{przyklad}[Podstawowe przestrzenie Frécheta]\label{ex:frechet-list11}
Każda przestrzeń Banacha jest przestrzenią Frécheta. Dalsze
przykłady to $\K^{\N}$ z topologią produktową oraz
$C^\infty([0,1])$ z seminormami
$p_n(f)=\max_{0\le j\le n}\sup_{t\in[0,1]}|f^{(j)}(t)|$,
gdzie pochodne na $(0,1)$ mają ciągłe przedłużenia na $[0,1]$.
W przestrzeni $\K^{\N}$ można wziąć seminormy
$p_n(x)=\max_{1\le j\le n}|x_j|$. Każda współrzędna ciągu
Cauchy'ego jest wtedy ciągiem Cauchy'ego w $\K$, a granice
współrzędnych dają granicę w topologii produktowej.
Wykażmy również zupełność $C^\infty([0,1])$. Jeśli $(f_n)$ jest ciągiem
Cauchy'ego we wszystkich seminormach $p_m$, to dla każdego $j\ge0$
ciąg $(f_n^{(j)})$ zbiega jednostajnie na $[0,1]$ do pewnej funkcji
ciągłej $g_j$. Z tożsamości
\[
 f_n^{(j)}(t)-f_n^{(j)}(0)
 =\int_{[0,t]}f_n^{(j+1)}(s)\,\mathrm d\lambda(s)
\]
i jednostajnej zbieżności funkcji podcałkowych otrzymujemy
$g_j(t)-g_j(0)=\int_{[0,t]}g_{j+1}(s)\,\mathrm d\lambda(s)$.
Twierdzenie podstawowe rachunku całkowego daje zatem
$g_j'=g_{j+1}$. Stąd $g_0\in C^\infty([0,1])$ oraz
$p_m(f_n-g_0)\to0$ dla każdego $m$.
\end{przyklad}

\subsection{Podprzestrzenie, iloczyny i przestrzenie ilorazowe}

\begin{twierdzenie}[Stabilność klasy przestrzeni Frécheta]\label{thm:frechet-stability11}
Niech $E$ będzie przestrzenią Frécheta, a $M\subseteq E$
domkniętą podprzestrzenią liniową. Wtedy $M$ oraz $E/M$
ze swoimi naturalnymi topologiami są przestrzeniami Frécheta.
Ponadto przeliczalny produkt przestrzeni Frécheta z topologią
produktową jest przestrzenią Frécheta.
\end{twierdzenie}
\begin{proof}
Niech $d$ będzie zgodną, zupełną i niezmienniczą względem
przesunięć metryką na $E$, jak w Stwierdzeniu~\ref{prop:frechet-metric11}.
Domknięta podprzestrzeń $M$ jest zupełna względem metryki
indukowanej z $d$, a lokalna wypukłość i metryzowalność są
dziedziczne dla podprzestrzeni.

Niech $q:E\to E/M$ będzie rzutowaniem ilorazowym.
Z domkniętości $M$ wynika, że
\[
 \bar d\bigl(q(x),q(y)\bigr)=\inf_{m\in M}d(x-y+m,0)
\]
jest metryką: nie zależy od przedstawicieli klas, spełnia
nierówność trójkąta, a $\bar d(q(x),q(y))=0$ oznacza
$x-y\in\overline M=M$. Ponadto
$q(B_d(x,r))=B_{\bar d}(q(x),r)$, więc $\bar d$ wyznacza
topologię ilorazową (rzutowanie $q$ jest otwarte).

Wykażemy zupełność. Niech $(u_n)$ będzie ciągiem Cauchy'ego
w $E/M$. Wybierzmy rosnący podciąg indeksów $(n_k)$ tak,
aby $\bar d(u_{n_{k+1}},u_{n_k})<2^{-k-1}$.
Wybierzmy przedstawiciela $z_1\in E$ klasy $u_{n_1}$.
Indukcyjnie wybieramy przedstawiciela $z_{k+1}$ klasy
$u_{n_{k+1}}$ z $d(z_{k+1},z_k)<2^{-k}$;
istnienie wynika z definicji $\bar d$. Ciąg $(z_k)$
jest Cauchy'ego w $E$, a więc zbiega do pewnego $z\in E$.
W konsekwencji $u_{n_k}=q(z_k)\to q(z)$.
Ponieważ całe $(u_n)$ jest ciągiem Cauchy'ego,
również $u_n\to q(z)$. Zatem $E/M$ jest zupełna.
Jej lokalna wypukłość wynika z obrazów otwartych,
wypukłych otoczeń zera przez $q$.

Niech teraz $(E_j)_{j\ge1}$ będzie ciągiem przestrzeni Frécheta.
Wybierzmy w każdym $E_j$ zgodną zupełną metrykę
$d_j$, niezmienniczą względem przesunięć. Metryka
\[
 D(x,y)=\sum_{j=1}^\infty 2^{-j}\min\{1,d_j(x_j,y_j)\}
\]
wyznacza topologię produktową na $\prod_{j\ge1}E_j$.
Ciąg $D$-Cauchy'ego jest Cauchy'ego na każdej współrzędnej;
po przejściu do granic współrzędnych i oszacowaniu ogona szeregu
otrzymujemy zbieżność względem $D$. Lokalna wypukłość
wynika z lokalnej wypukłości czynników.
\end{proof}

\begin{uwaga}[Znaczenie domkniętości podprzestrzeni]
Jeśli $M$ nie jest domknięta, iloraz $E/M$ z topologią
ilorazową nie jest Hausdorffa. Dlatego nie jest przestrzenią
Frécheta. Warunek domkniętości w Twierdzeniu~\ref{thm:frechet-stability11}
nie może zostać pominięty.
\end{uwaga}

\subsection{Twierdzenie o odwzorowaniu otwartym}

\begin{lemat}[Krok Baire'a]\label{lem:baire-open11}
Niech $E$ będzie TVS, $F$ lokalnie wypukłą TVS Hausdorffa
będącą przestrzenią Baire'a, a $T:E\to F$ ciągłym,
surjektywnym operatorem liniowym. Jeśli $U\subseteq E$
jest otwartym, absolutnie wypukłym otoczeniem zera, to
$\overline{T(U)}$ jest otoczeniem zera w $F$.
\end{lemat}
\begin{proof}
Ponieważ $U$ jest pochłaniający, a $T$ jest surjektywny,
$F=\bigcup_{n=1}^\infty nT(U)$. Zbiór $nT(U)$ jest
wypukły i zrównoważony, a jego domknięcie również.
Z własności Baire'a wynika, że pewien zbiór
$\overline{nT(U)}=n\overline{T(U)}$ ma niepuste wnętrze.
Stąd także $\overline{T(U)}$ ma niepuste wnętrze.
Jeżeli wypukły, zrównoważony zbiór $C$ ma niepuste wnętrze,
to zawiera otoczenie zera: dla $a\in\operatorname{int}C$
mamy także $-a\in\operatorname{int}C$, a wypukłość
daje otoczenie zera zawarte w $C$, jako zbiór
$\tfrac12((a+V)+(-a+V))=V$ dla dostatecznie małego,
otwartego, zrównoważonego otoczenia $V$ zera.
Stosujemy to do $C=\overline{T(U)}$.
\end{proof}

\begin{twierdzenie}[Twierdzenie o odwzorowaniu otwartym]\label{thm:open-mapping11}
Niech $E$ i $F$ będą przestrzeniami Frécheta nad $\K$.
Jeżeli $T:E\to F$ jest ciągłym i surjektywnym
operatorem liniowym, to $T$ jest odwzorowaniem otwartym:
dla każdego otwartego $O\subseteq E$ zbiór $T(O)$
jest otwarty w $F$.
\end{twierdzenie}
\begin{proof}
Wystarczy pokazać, że obraz dowolnego otoczenia zera
w $E$ zawiera otoczenie zera w $F$. Ustalmy takie
otoczenie $U$ i zmniejszmy je, jeśli trzeba, aby było
otwarte i absolutnie wypukłe. Wybierzmy otwarte,
absolutnie wypukłe otoczenie $U_0$ zera takie, że
$\overline{U_0}\subseteq U$. Istnienie wynika z lokalnej
wypukłości i ciągłości seminorm wyznaczających topologię.

Indukcyjnie konstruujemy otwarte, absolutnie wypukłe
otoczenia $(U_n)_{n\ge1}$ zera w $E$, takie że
$U_n+U_n\subseteq U_{n-1}$ oraz $U_n\subseteq A_n$,
gdzie $(A_n)$ jest ustaloną malejącą przeliczalną bazą
otoczeń zera w $E$. Możemy to zrobić, zmniejszając
kolejne otoczenia. Wówczas $(U_n)$ jest bazą otoczeń zera.

Z Lematu~\ref{lem:baire-open11}, zastosowanego osobno
do każdego $U_n$, wynika, że
$\overline{T(U_n)}$ zawiera otoczenie zera w $F$.
Wybieramy zatem otwarte, absolutnie wypukłe otoczenia
$V_n$ zera w $F$ takie, że
$V_n\subseteq\overline{T(U_n)}\cap B_n$, gdzie $(B_n)$
jest malejącą przeliczalną bazą otoczeń zera w $F$.

Niech $y\in V_1$. Wobec $y\in\overline{T(U_1)}$
i otwartości $V_2$ istnieje $x_1\in U_1$ takie, że
$y-Tx_1\in V_2$. Ponieważ $V_2\subseteq\overline{T(U_2)}$,
możemy wybrać $x_2\in U_2$ tak, aby
$y-T(x_1+x_2)\in V_3$. Kontynuując konstrukcję,
otrzymujemy $x_n\in U_n$ spełniające
\[
 r_n:=y-T\Bigl(\sum_{k=1}^n x_k\Bigr)\in V_{n+1}
 \quad (n\ge1).
\]
Ponieważ $V_{n+1}\subseteq B_{n+1}$, mamy $r_n\to0$ w $F$.

Pokażemy, że szereg $\sum_{n=1}^\infty x_n$ zbiega w $E$.
Z inkluzji $U_n+U_n\subseteq U_{n-1}$ wynika także
$U_n\subseteq U_{n-1}$. Indukcja wsteczna dowodzi, że
$\sum_{k=m}^N x_k\in U_{m-1}$ dla dowolnych $N\ge m\ge1$.
Ponieważ $(U_n)$ jest bazą otoczeń zera, ciąg sum częściowych
jest Cauchy'ego. Zupełność $E$ daje
$x=\lim_{N\to\infty}\sum_{k=1}^N x_k$.
Dla $m=1$ sumy częściowe należą do $U_0$, więc
$x\in\overline{U_0}\subseteq U$. Z ciągłości $T$ i
$r_n\to0$ otrzymujemy $Tx=y$. Udowodniliśmy
$V_1\subseteq T(U)$.

Dla dowolnego zbioru otwartego $O\subseteq E$ i punktu
$y=Tx\in T(O)$ wybieramy otoczenie zera $U_x$ z
$x+U_x\subseteq O$. Wówczas $T(O)$ zawiera otoczenie
$y$ postaci $y+V$, gdzie $V\subseteq T(U_x)$.
Zatem $T(O)$ jest otwarty.
\end{proof}

\begin{uwaga}[Co jest rzeczywiście potrzebne?]\label{rem:open-assumptions11}
Dowód wykorzystuje zupełność i metryzowalność dziedziny,
własność Baire'a oraz metryzowalność przestrzeni wartości,
a także lokalną wypukłość obu przestrzeni. W konsekwencji
ten sam argument działa dla $E$ będącej przestrzenią Frécheta
oraz $F$ będącej metryzowalną lokalnie wypukłą TVS Hausdorffa
o własności Baire'a, nawet jeśli zupełności $F$ nie założono.
Nie twierdzimy natomiast, że dowolny ciągły operator
liniowy między dwiema przestrzeniami Frécheta jest otwarty:
konieczna jest surjektywność.
\end{uwaga}

\begin{wniosek}[Twierdzenie o ciągłości odwrotności]\label{cor:inverse-mapping11}
Niech $E,F$ będą przestrzeniami Frécheta, a
$T\in\mathcal L(E,F)$ liniową bijekcją.
Wtedy $T^{-1}:F\to E$ jest ciągłym operatorem liniowym.
Zatem $T$ jest izomorfizmem liniowo-topologicznym.
\end{wniosek}
\begin{proof}
Z Twierdzenia~\ref{thm:open-mapping11} operator $T$ jest otwarty.
Dla każdego otwartego $U\subseteq E$ zbiór
$(T^{-1})^{-1}(U)=T(U)$ jest zatem otwarty w $F$.
\end{proof}

\begin{wniosek}[Oszacowanie normy operatora odwrotnego]\label{cor:bounded-inverse11}
Niech $E,F$ będą przestrzeniami Banacha i niech
$T:E\to F$ będzie liniowym, ciągłym i bijektywnym
operatorem. Istnieje stała $C>0$ taka, że
$\|T^{-1}y\|_E\le C\|y\|_F$ dla każdego $y\in F$.
Równoważnie, $\|Tx\|_F\ge C^{-1}\|x\|_E$ dla $x\in E$.
\end{wniosek}
\begin{proof}
Z Wniosku~\ref{cor:inverse-mapping11} odwrotność jest ciągła.
Liniowy operator między przestrzeniami unormowanymi jest
ciągły wtedy i tylko wtedy, gdy jest ograniczony; przyjmujemy
$C=\max\{1,\|T^{-1}\|\}$.
\end{proof}

\begin{wniosek}[Izomorfizm z przestrzenią ilorazową]\label{cor:quotient-open11}
Niech $E,F$ będą przestrzeniami Frécheta oraz
$T\in\mathcal L(E,F)$ operatorem surjektywnym.
Wtedy $M=\ker T$ jest domkniętą podprzestrzenią $E$,
a operator
$\widetilde T:E/M\to F$, $\widetilde T(x+M)=Tx$,
jest izomorfizmem liniowo-topologicznym.
\end{wniosek}
\begin{proof}
Ciągłość $T$ i domkniętość $\{0\}\subset F$ dają
domkniętość $M$. Własność uniwersalna ilorazu zapewnia
ciągłość $\widetilde T$. Jest ona liniową bijekcją.
Ponieważ $E/M$ i $F$ są przestrzeniami Frécheta,
Wniosek~\ref{cor:inverse-mapping11} daje ciągłość odwrotności.
\end{proof}

\subsection{Twierdzenie o wykresie domkniętym}

\begin{definicja}[Wykres operatora]\label{def:graph11}
Jeżeli $E,F$ są TVS, a $T:D(T)\subseteq E\to F$
odwzorowaniem, to jego \emph{wykresem} nazywamy
$G(T)=\{(x,Tx):x\in D(T)\}\subseteq E\times F$.
Dla operatora liniowego $G(T)$ jest podprzestrzenią
liniową produktu.
\end{definicja}

\begin{stwierdzenie}[Kryterium ciągowe domknięcia wykresu]
\label{prop:graph-sequential11}
Niech $E,F$ będą metryzowalnymi TVS Hausdorffa, a
$T:E\to F$ operatorem liniowym określonym na całym $E$.
Wykres $G(T)$ jest domknięty w $E\times F$ wtedy i tylko
wtedy, gdy dla każdego ciągu $(x_n)$ w $E$ i każdego
$x\in E$, $y\in F$ spełniających $x_n\to x$ i
$Tx_n\to y$ zachodzi $Tx=y$.
\end{stwierdzenie}
\begin{proof}
Produkt $E\times F$ jest metryzowalny. W przestrzeni
metrycznej zbiór jest domknięty dokładnie wtedy, gdy
zawiera granicę każdego zbieżnego ciągu swoich punktów.
Zastosowanie do ciągu $(x_n,Tx_n)$ daje tezę.
\end{proof}

\begin{twierdzenie}[Twierdzenie o wykresie domkniętym]\label{thm:closed-graph11}
Niech $E,F$ będą przestrzeniami Frécheta, a
$T:E\to F$ operatorem liniowym określonym na całym $E$.
Wówczas $T$ jest ciągły wtedy i tylko wtedy,
gdy jego wykres $G(T)$ jest domknięty w $E\times F$.
\end{twierdzenie}
\begin{proof}
Jeżeli $T$ jest ciągły, to odwzorowanie
$(x,y)\mapsto y-Tx$ jest ciągłe z $E\times F$ do $F$.
Jego przeciwobraz zbioru domkniętego $\{0\}$ to $G(T)$,
a więc wykres jest domknięty.

Odwrotnie, załóżmy, że $G(T)$ jest domknięty.
Iloczyn $E\times F$ jest przestrzenią Frécheta
na mocy Twierdzenia~\ref{thm:frechet-stability11},
a jego domknięta podprzestrzeń $G(T)$ również.
Rzutowanie
$P:G(T)\to E$, $P(x,Tx)=x$, jest ciągłą liniową
bijekcją między przestrzeniami Frécheta. Z Wniosku
\ref{cor:inverse-mapping11} wynika ciągłość $P^{-1}$.
Stąd $T=\pi_F\circ P^{-1}$ jest ciągły, gdzie
$\pi_F:E\times F\to F$ oznacza drugie rzutowanie.
\end{proof}

\begin{center}
\begin{tikzpicture}[>=Stealth, every node/.style={font=\small}]
  \node (g) at (0,1.65) {$G(T)\subseteq E\times F$};
  \node (e) at (-2.75,0) {$E$};
  \node (f) at (2.75,0) {$F$};
  \draw[->] (g) -- node[above left] {$P$} (e);
  \draw[->] (g) -- node[above right] {$\pi_F|_{G(T)}$} (f);
  \draw[->,dashed] (e) -- node[below] {$T=\pi_F\circ P^{-1}$} (f);
\end{tikzpicture}

\smallskip
\small Schemat dowodu twierdzenia o wykresie domkniętym:
$P$ jest ciągłą bijekcją, a jego odwrotność jest ciągła
na mocy twierdzenia o odwzorowaniu otwartym.
\end{center}

\begin{uwaga}[Pełna dziedzina operatora jest niezbędna]
Twierdzenie~\ref{thm:closed-graph11} wymaga, by $T$ był
określony na całym $E$. Operatory różniczkowe jako
operatory nieograniczone w przestrzeniach Banacha mają
zwykle właściwe dziedziny $D(T)\subsetneq E$ i mogą
mieć domknięte wykresy bez ciągłości w normie $E$.
Przykład podamy poniżej.
\end{uwaga}

\subsection{Kontrprzykłady i zastosowania}

\begin{przyklad}[Brak zupełności przeciwdziedziny: bijekcja nieotwarta]
\label{ex:incomplete-codomain11}
Na zbiorze $\ell^1$ rozważmy dwie normy
$\|x\|_1=\sum_{n\ge1}|x_n|$ oraz
$\|x\|_2=(\sum_{n\ge1}|x_n|^2)^{1/2}$.
Operator identycznościowy
$J:(\ell^1,\|\cdot\|_1)\to(\ell^1,\|\cdot\|_2)$
jest ciągłą liniową bijekcją, bo $\|x\|_2\le\|x\|_1$.
Jego odwrotność nie jest ciągła. Dla wektora $u^{(n)}$
o pierwszych $n$ współrzędnych równych $1/n$
i pozostałych równych zeru mamy
$\|u^{(n)}\|_1=1$ i $\|u^{(n)}\|_2=n^{-1/2}\to0$.
Zatem $J$ nie jest otwarty. Przeciwdziedzina
$(\ell^1,\|\cdot\|_2)$ nie jest zupełna: jej
uzupełnieniem jest $\ell^2$, ponieważ $c_{00}$ jest
gęste w $\ell^2$.
\end{przyklad}

\begin{przyklad}[Wykres domknięty nie zapewnia ciągłości bez zupełności]
\label{ex:incomplete-domain-graph11}
Niech $E=(\ell^1,\|\cdot\|_2)$, $F=(\ell^1,\|\cdot\|_1)$,
a $T:E\to F$ będzie odwzorowaniem identycznościowym.
Nie jest ono ciągłe, jak pokazuje ciąg $u^{(n)}$ z
Przykładu~\ref{ex:incomplete-codomain11}. Ma jednak
wykres domknięty w $E\times F$. Istotnie, jeśli
$x_n\to x$ w normie $\ell^2$ oraz
$Tx_n=x_n\to y$ w normie $\ell^1$, to $x_n\to y$
także w normie $\ell^2$. Jednoznaczność granicy
w $\ell^2$ daje $x=y$, czyli $(x,y)\in G(T)$.
Kryterium ciągowe domknięcia wykresu kończy dowód.
Tutaj $E$ jest niezupełna, mimo że $F$ jest przestrzenią Banacha.
\end{przyklad}

\begin{przyklad}[Domknięty nieograniczony operator różniczkowania]
\label{ex:closed-derivative11}
Niech $E=F=C([0,1])$ z normą supremum i
$D(T)=C^1([0,1])$. Określmy $Tf=f'$. Jest to liniowy
operator nieograniczony względem normy $E$:
$f_n(t)=n^{-1}\sin(nt)$ spełnia $\|f_n\|_\infty\to0$,
podczas gdy $\|Tf_n\|_\infty=1$.
Wykres $G(T)$ jest jednak domknięty w $E\times F$.
Jeżeli $f_n\to f$ i $f_n'\to g$ jednostajnie,
to dla $t\in[0,1]$ z twierdzenia podstawowego rachunku
całkowego mamy
\[
 f_n(t)-f_n(0)=\int_{[0,t]}f_n'(s)\,\mathrm d\lambda(s).
\]
Przechodząc do granicy, otrzymujemy
$f(t)-f(0)=\int_{[0,t]}g(s)\,\mathrm d\lambda(s)$.
Ponieważ $g$ jest ciągła, funkcja $f$ jest klasy $C^1$
i $f'=g$. Nie przeczy to Twierdzeniu~\ref{thm:closed-graph11},
ponieważ $D(T)\ne E$.
\end{przyklad}

\begin{przyklad}[Operator pochodnej w przestrzeni Frécheta]
\label{ex:frechet-derivative11}
Niech $E=F=C^\infty([0,1])$ z seminormami $p_n$
z Przykładu~\ref{ex:frechet-list11} i niech $Df=f'$.
Operator $D:E\to F$ jest ciągły, ponieważ
$p_n(Df)\le p_{n+1}(f)$ dla każdego $n$.
Jest również surjektywny: dla każdego $g\in F$ funkcja
$Rg(t)=\int_{[0,t]}g(s)\,\mathrm d\lambda(s)$ należy do $E$
i spełnia $D(Rg)=g$. Jądro $D$ składa się z funkcji stałych.
Zatem Twierdzenie~\ref{cor:quotient-open11} daje
izomorfizm liniowo-topologiczny
$C^\infty([0,1])/\K\cong C^\infty([0,1])$,
gdzie $\K$ oznacza podprzestrzeń funkcji stałych.
W istocie $R$ jest ciągły, ponieważ
$p_n(Rg)\le 2p_n(g)$ dla wszystkich $n\ge0$.
\end{przyklad}

\begin{przyklad}[Gęsty, niedomknięty obraz operatora]
\label{ex:dense-range11}
Naturalne włożenie $\iota:\ell^1\to\ell^2$ jest
ciągłe i różnowartościowe. Jego obraz jest gęsty, gdyż
zawiera $c_{00}$, ale nie jest cały: ciąg
$(n^{-3/4})_{n\ge1}$ należy do $\ell^2\setminus\ell^1$.
Odwrotność $\iota^{-1}:\iota(\ell^1)\to\ell^1$
nie jest ciągła przy topologii indukowanej z $\ell^2$.
Twierdzenie o odwzorowaniu otwartym nie ma zastosowania
do operatora $\ell^1\to\iota(\ell^1)$, ponieważ
przeciwdziedzina z topologią indukowaną nie jest zupełna.
\end{przyklad}

\begin{wniosek}[Równoważność zupełnych norm porównywalnych]
\label{cor:two-banach-norms11}
Niech $V$ będzie przestrzenią wektorową nad $\K$,
a $\|\cdot\|_a$ oraz $\|\cdot\|_b$ normami na $V$
takimi, że obie przestrzenie $(V,\|\cdot\|_a)$ i
$(V,\|\cdot\|_b)$ są Banacha. Jeśli istnieje $C>0$
z $\|x\|_b\le C\|x\|_a$ dla wszystkich $x\in V$,
to istnieje $D>0$ takie, że
$\|x\|_a\le D\|x\|_b$ dla wszystkich $x\in V$.
\end{wniosek}
\begin{proof}
Operator identycznościowy między podanymi przestrzeniami
Banacha jest ciągłą bijekcją. Jego odwrotność jest ciągła
na mocy Wniosku~\ref{cor:inverse-mapping11}, co jest
równoważne żądanemu oszacowaniu.
\end{proof}

\begin{uwaga}[Pomost do rachunku różniczkowego]
W dalszych rozdziałach będziemy badać pochodne Gâteaux
i Frécheta odwzorowań między przestrzeniami Banacha.
W twierdzeniu o lokalnej funkcji odwrotnej zakłada się
odwracalność pochodnej Frécheta, która jest ograniczonym
operatorem liniowym. Wniosek~\ref{cor:bounded-inverse11}
pokazuje, dlaczego algebraiczna odwracalność takiej
pochodnej między przestrzeniami Banacha zapewnia
ciągłość jej odwrotności. Nie należy mylić tego z
odwracalnością samego odwzorowania nieliniowego.
\end{uwaga}
